\section{Testy poprawności implementacji}

Przeprowadzono następujące testy poprawności implementacji:

\subsection{Weryfikacja modelu jednokryterialnego}
\begin{enumerate}
  \item \textbf{Test ograniczeń pojemności maszyn} - sprawdzono, czy dla każdego miesiąca i typu maszyny całkowity czas produkcji nie przekracza dostępnego czasu.
  \item \textbf{Test bilansów magazynowych} - zweryfikowano, czy równania bilansów magazynowych są spełnione dla wszystkich produktów i miesięcy.
  \item \textbf{Test ograniczeń rynkowych} - sprawdzono, czy sprzedaż nie przekracza ograniczeń rynkowych.
  \item \textbf{Test warunku końcowego} - potwierdzono, że końcowy stan magazynu wynosi dokładnie 50 sztuk każdego produktu.
\end{enumerate}

\subsection{Weryfikacja modelu dwukryterialnego}
\begin{enumerate}
  \item \textbf{Test generowania scenariuszy} - sprawdzono, czy wygenerowane scenariusze dochodów mają wartości w zakresie [5; 12].
  \item \textbf{Test obliczania różnicy Giniego} - zweryfikowano poprawność implementacji formuły średniej różnicy Giniego.
  \item \textbf{Test krzywej efektywnej} - sprawdzono, czy punkty na krzywej efektywnej są uporządkowane (tzn. czy większemu zyskowi odpowiada większe ryzyko).
  \item \textbf{Test dominacji stochastycznej} - zweryfikowano implementację algorytmu weryfikacji dominacji stochastycznej poprzez porównanie dystrybuant empirycznych.
\end{enumerate}

\subsection{Wyniki testów}

Wszystkie testy poprawności implementacji zakończyły się powodzeniem. Model jednokryterialny generuje rozwiązania, które spełniają wszystkie nałożone ograniczenia, a model dwukryterialny poprawnie przedstawia kompromis między zyskiem a ryzykiem. Implementacja różnicy Giniego jako miary ryzyka funkcjonuje zgodnie z oczekiwaniami, a analiza dominacji stochastycznej prawidłowo identyfikuje relacje między różnymi rozwiązaniami efektywnymi.

